% ---------------------------------------------------------------------------- %
% Pacotes
\usepackage[T1]{fontenc}
\usepackage[brazil]{babel}           % define a linguagem padrão
\usepackage[utf8x]{inputenc}
\usepackage[pdftex]{graphicx}           % usamos arquivos pdf/png como figuras
\usepackage{setspace}                   % espaçamento flexível
\usepackage{makeidx}                    % índice remissivo
\usepackage[nottoc]{tocbibind}          % acrescentamos a bibliografia/indice/conteudo no Table of Contents
\usepackage{courier}                    % usa o Adobe Courier no lugar de Computer Modern Typewriter
\usepackage{type1cm}                    % fontes realmente escaláveis
\usepackage{listings}                   % para formatar código-fonte (ex. em Java)
\usepackage{titletoc}
\usepackage[fixlanguage]{babelbib}
\usepackage[font=small,format=plain,labelfont=bf,up,textfont=it,up]{caption}
\usepackage[usenames,svgnames,dvipsnames]{xcolor}
\usepackage[a4paper,top=2.54cm,bottom=2.0cm,left=2.0cm,right=2.54cm]{geometry} % margens
\usepackage[pdftex,plainpages=false,pdfpagelabels,pagebackref,colorlinks=true,citecolor=black,linkcolor=black,urlcolor=black,filecolor=black,bookmarksopen=true]{hyperref} % links em preto
\usepackage[all]{hypcap}                % soluciona o problema com o hyperref e capitulos
\usepackage{indentfirst}                % recuo também no primeiro parágrafo após título (norma ABNT)

\fontsize{60}{62}\usefont{OT1}{cmr}{m}{n}{\selectfont}
\usepackage{color} % colorir texto
\usepackage{xcolor}
\usepackage{amsmath} %expressoes matematicas
\usepackage{pdflscape} %permite mudar orientacao da pagina para paisagem E MUDA NO PDF
\usepackage{longtable} %permite tabelas que ocupem mais de uma página
\usepackage{multicol}
\usepackage{enumitem} %para pode iniciar enumerate em 0
\usepackage{pdfpages}


% ---------------------------------------------------------------------------- %
% Cabeçalhos similares ao TAOCP de Donald E. Knuth
\usepackage{fancyhdr}
\pagestyle{fancy}
\fancyhf{}
\renewcommand{\chaptermark}[1]{\markboth{\MakeUppercase{#1}}{}}
\renewcommand{\sectionmark}[1]{\markright{\MakeUppercase{#1}}{}}
\renewcommand{\headrulewidth}{0pt}

% ---------------------------------------------------------------------------- %
% Estilo de capítulo compacto (sem o espaço em branco padrão do LaTeX antes do título)
\usepackage{titlesec}
\titleformat{\chapter}[hang]{\normalfont\bfseries\LARGE}{\thechapter}{1em}{\MakeUppercase}
\titleformat{name=\chapter,numberless}[block]{\normalfont\bfseries\LARGE\centering}{}{0pt}{}
\titlespacing*{\chapter}{0pt}{0pt}{40pt}

% ---------------------------------------------------------------------------- %
\graphicspath{{./figuras/}}             % caminho das figuras (recomendável)
\frenchspacing                          % arruma o espaço: id est (i.e.) e exempli gratia (e.g.)
\urlstyle{same}                         % URL com o mesmo estilo do texto e não mono-spaced
\makeindex                              % para o índice remissivo
\raggedbottom                           % para não permitir espaços extra no texto
\fontsize{60}{62}\usefont{OT1}{cmr}{m}{n}{\selectfont}
\normalsize

% ---------------------------------------------------------------------------- %
% Opções de listing usados para o código fonte
% Ref: http://en.wikibooks.org/wiki/LaTeX/Packages/Listings
\lstset{ %
language=Java,                  % choose the language of the code
basicstyle=\footnotesize,       % the size of the fonts that are used for the code
numbers=left,                   % where to put the line-numbers
numberstyle=\footnotesize,      % the size of the fonts that are used for the line-numbers
stepnumber=1,                   % the step between two line-numbers. If it's 1 each line will be numbered
numbersep=5pt,                  % how far the line-numbers are from the code
showspaces=false,               % show spaces adding particular underscores
showstringspaces=false,         % underline spaces within strings
showtabs=false,                 % show tabs within strings adding particular underscores
frame=single,	                % adds a frame around the code
framerule=0.6pt,
tabsize=2,	                    % sets default tabsize to 2 spaces
captionpos=b,                   % sets the caption-position to bottom
breaklines=true,                % sets automatic line breaking
breakatwhitespace=false,        % sets if automatic breaks should only happen at whitespace
escapeinside={\%*}{*)},         % if you want to add a comment within your code
backgroundcolor=\color[rgb]{1.0,1.0,1.0}, % choose the background color.
rulecolor=\color[rgb]{0.8,0.8,0.8},
extendedchars=true,
xleftmargin=10pt,
xrightmargin=10pt,
framexleftmargin=10pt,
framexrightmargin=10pt
}

\lstdefinelanguage{blah} % minha linguagem se chama blah
{
% Aqui vão os atributos dela. Para ver mais atributos, Google!
tabsize=2, % o tab possui 2 espaços
frame=shadowbox,% o tipo de box que usaremos
%rulesepcolor=\color{lightgrey}, %/ a cor pra fazer o efeito legal da box
captionpos=b,% o caption do código vai para baixo (botton) o default eh top
extendedchars=false, % eu usei isso para conseguir colocar expressoes matematicas utilizando o scapeinside tbm
escapeinside='', % usando aspas, pode-se usar qualquer comando dentro deles
morekeywords = {if, else, while, do, for, and, or, action, initialize} % palavras reservadas que ficaram em negrito
}


% ABNT %%%%%%%%%%%%%%%%%%%%%%%%%%%%%%%%%%%%%%%%%%%%%%%%%%%%%%%%%%%%%%%%%%%%%%%%%%%%%
\usepackage[alf,% Para usar o estilo de citacao alfabetico (seção 9.2 em 6023/2000
abnt-etal-text=emph,% o et al. vai aparecer em italico
abnt-and-type=e,% separador de autores com ´e´
bibjustif,% bibliografia justificada no final
abnt-etal-cite=3]{abntex2cite}
% No caso, a 6023 diz que se deve usar ‘et al.’ para mais de três autores, ou seja, a partir de quatro autores. Então, usarei abnt-etal-cite=3
\usepackage{cite}%%% NOVO

% ------------ colocar o [fulanclearo 2002] antes da descricao de referencia bibliografica
\makeatletter
\renewcommand\@biblabel[1]{[#1]~}% colocar citação na frente da lista de bibliografia
\renewcommand\citepunct{;~}%%% NOVO
%\renewcommand{\@biblabel}[1]{[\textbf{#1}]~}
\makeatother
\citebrackets[] % fazer com que as citacoes dentro do texto virem colchetes
% -------------------------------------------------------
%%%%%%%%%%%%%%%%%%%%%%%%%%%%%%%%%%%%%%%%%%%%%%%%%%%%%%%%%%%%%%%%%%%%%%%%%%%%%%%%%%%%
